
Реальный платёжный продукт сопротивляется после запуска. Тогда
выясняется, что корректно спроектированный поток приносит
убыток, теряет деньги на сверке, упирается в~комплаенс, провоцирует
ложные отказы, накапливает спорные списания и~ломается там, где на
уровне протокола всё выглядело безупречно. К~моменту входа в~эту
часть у~читателя есть соблазн думать о~платёжной системе как
о~красивой архитектуре: роли определены, протоколы разобраны,
различия между сетями и системами очерчены. Эта картина не выдерживает
прода.

Часть~IV переводит разговор с~языка \enquote{как должен работать
платёж} на язык того, что происходит с~платежом уже в~проде.
Здесь недостаточно знать, что такое
авторизация, токен или перенос ответственности. Нужно видеть, как эти
механизмы ведут себя под давлением реальных атак, операционных ошибок,
требований регулятора, продуктовых компромиссов и поведения
пользователя.

Регуляторный блок открывает часть, потому что без него разговор о~фроде,
спорах и~сверке лишён опоры. Сначала идёт регуляторика:
161-ФЗ, положения Банка России, лицензионные модели и~сама форма,
в~которой норма попадает в~инженерный план. Затем разбираются
дисциплины комплаенса, которые в~реальной системе сшиты вместе,
но требуют разного подхода: идентификация клиентов и~ТСП
(в~международной терминологии~-- KYC, Know Your Customer,
и~KYB, Know Your Business), ПОД/ФТ
(противодействие легализации\,/\,отмыванию доходов и финансированию
терроризма; в~международной практике~-- AML, anti-money
laundering) и~санкционный скрининг. Они разнесены по отдельным главам сознательно: пороги,
перечни и~маршруты эскалации у~каждой дисциплины свои, а~их смешение
быстро оборачивается ошибками на стыке.

Дальше~--- прямое давление со стороны среды. Фрод~--- набор конкретных атак,
вынуждающих балансировать между конверсией и~защитой. Подписки
и~Card-on-File показывают, что у~платёжного инструмента долгая жизнь.
Даже когда клиент изначально согласен платить, эта
жизнь требует дисциплины, не сводимой к~одной успешно
прошедшей транзакции.

Заключительная половина части посвящена последствиям уже совершившихся
операций и~цене выбора платёжного метода. Сверка отвечает на
вопрос, где именно теряются деньги или состояния, когда внешние
события перестают совпадать с~внутренним учётом. Диспуты
и~процессы возвратных платежей показывают, что успешный платёж
не равен защищённой выручке. Управление отказами показывает, как сеть, эмитент,
эквайер, антифрод и пользовательская логика сталкиваются в~одной
точке.

Часть собирает вместе темы, которые в~компаниях разнесены по разным
командам. Регуляторика и~ПОД/ФТ живут у~комплаенса, фрод~--- у~команды
риска, регулярные списания~--- у~биллинга, сверка~--- у~финансовых операций,
возвратные платежи~--- у~команды диспутов, управление отказами~--- у~команды
роста или платёжной эффективности. Для платёжного продукта это одна
и~та же реальность. Пользователь не видит внутренних разделений между
подразделениями; он видит, что оплата прошла, не прошла,
была оспорена, повторилась неожиданно или исчезла где-то между
авторизацией и выпиской.

Продукт начинает взрослеть не в~тот момент, когда он впервые принял платёж,
а~в~тот, когда команда научилась связывать конверсию с~потерями
от фрода, рост~--- с~требованиями ПОД/ФТ, доход от подписок~--- с~дисциплиной
согласия, а~финансовый результат~--- со сверкой, диспутами и отказами.
Тут разговор смещается с~платёжной инфраструктуры
к~платёжной эксплуатации.

Дальше в~части~\ref{part:architecture} книга переходит к~внутренней
топологии платёжной системы: какой компонент за что отвечает
и~как они связаны.
